\documentclass[10pt,UTF8]{article}


\usepackage{geometry}

\usepackage{ctex}

\geometry{a4paper,scale=0.9}
\ctexset{today=old}

% \usepackage{times}  % 设置正文字体为 Adobe Times Roman
% \usepackage{mathptmx} % 只加载数学字体
% \renewcommand{\rmdefault}{cmr} % 恢复正文字体为 Computer Modern Roman

\usepackage{bm} % 数学粗体（避免斜体被取消）

% \usepackage{makecell}
% \usepackage{multirow} % 表格多行合并

% \usepackage{graphicx} %插入图片的宏包
% \usepackage{subfigure}
% \usepackage{float} %设置图片浮动位置的宏包

\usepackage{amsmath}
\usepackage{esint} % 提供更丰富的积分符号
\usepackage{amssymb}
\usepackage{mathtools}

\usepackage{physics}

% \usepackage{siunitx}
% % 关闭警告
% \ExplSyntaxOn
% \msg_redirect_name:nnn { siunitx } { physics-pkg } { none }
% \ExplSyntaxOff

% \usepackage{bussproofs} % 自然演绎推理树

% 交叉引用
% 关闭边框，设置颜色
\usepackage[colorlinks=true, linkcolor=darkgray, citecolor=teal, urlcolor=blue]{hyperref}
% 使 \cref 和 \Cref 可用
\usepackage{cleveref}

\usepackage{bookmark}  % 生成 PDF 大纲

\allowdisplaybreaks[3]  % 允许换页


%%%%%%%%%%%%%%%%%%%%%%%% tcolorbox %%%%%%%%%%%%%%%%%%%%%%%%%%
\usepackage[most]{tcolorbox}

\newenvironment{tipBox}
{\begin{tcolorbox}[colback=green!5!white,colframe=green!75!black,parbox=false,before upper=\par,before lower=\par]}
{\end{tcolorbox}}

\newenvironment{conceptBox}
{\begin{tcolorbox}[colback=blue!5!white,colframe=blue!75!black]}
{\end{tcolorbox}}

\newenvironment{theoremBoxNoIndent}
{\begin{tcolorbox}[colback=yellow!5!white,colframe=yellow!75!black]}
{\end{tcolorbox}}

\newenvironment{definitionBox}
{\begin{tcolorbox}[colback=blue!5!white,colframe=blue!75!black,parbox=false,before upper=\par,before lower=\par]}
{\end{tcolorbox}}

\newenvironment{theoremBox}
{\begin{tcolorbox}[colback=yellow!5!white,colframe=yellow!75!black,parbox=false,before upper=\par,before lower=\par]}
{\end{tcolorbox}}
%%%%%%%%%%%%%%%%%%%%%%%%%%%%%%%%%%%%%%%%%%%%%%%%%%%%%%%%%%%%
            
%%%%%%%%%%%%%%%%%%%%% amsthm %%%%%%%%%%%%%%%%
\usepackage{amsthm}

\theoremstyle{definition}
\newtheorem{definition inner}{Definition}[section]
\newenvironment{definition}[1][]
{\begin{definitionBox}\begin{definition inner}[#1]\hfill\par}
{\end{definition inner}\end{definitionBox}}

\theoremstyle{definition}
\newtheorem{theorem inner}{Theorem}[section]
\newenvironment{theorem}[1][]
{\begin{theoremBox}\begin{theorem inner}[#1]\hfill\par}
{\end{theorem inner}\end{theoremBox}}

\theoremstyle{definition}
\newtheorem*{proof inner}{\textit{Proof}}
\renewenvironment{proof}[1][]
{\begin{proof inner}[#1]\hfill\par}
{\qed\end{proof inner}}

\theoremstyle{plain}
\newtheorem{tip inner}{Tip}[section]
\newenvironment{tip}[1][]
{\begin{tipBox}\begin{tip inner}[#1]\hfill\par}
{\end{tip inner}\end{tipBox}}

\theoremstyle{remark}
\newtheorem*{remark}{\textit{Remark}}
%%%%%%%%%%%%%%%%%%%%%%%%%%%%%%%%%%%%%%%%%%%%%%


\newcommand{\iu}{\mathrm{i}}  % 虚数单位 i


\begin{document}

\part*{First Order ODEs}

The \textbf{general form} of a first order ODE is
\[
    \dv{y}{t} = f(t,y)
\]
for some given function \(f\).

\section{Linear}

The \textbf{general form} of a first order linear ODE is
\[
    \dv{y}{t} + p(t)\cdot y = q(t)
\]
for some given functions \(p\) and \(q\).

\subsection{Method of Integrating Factors}

\begin{theoremBox}
    For ODEs of the form
    \[
        \dv{y}{t} + p(t)\cdot y = q(t)
    \]
    the \textbf{general solution} is
    \[
        y = \frac{1}{\mu(t)} \int \mu(t)\cdot q(t)\dd{t}
    \]
    where the \textbf{integrating factor} is any \(\mu(t)\) that satisfies
    \[
        \ln\abs{\mu(t)} = \int p(t)\dd{t}
    \]
\end{theoremBox}
\begin{remark}
    Indefinite integrals yield a family of functions differing by a constant.
    Each indefinite integral results in an independent constant, which represents a degree of freedom.
    However, the constant from \(\mu(t)\) cancels out in the final solution for \(y\).
\end{remark}


\subsection{Variation of Parameters}

\begin{theoremBox}
    For ODEs of the form
    \[
        \dv{y}{t} + p(t)\cdot y = q(t)
    \]
    let \(h(t)\) be a function such that
    \[
        h(t) = \int p(t)\dd{t}
    \]
    Then the \textbf{general solution} is given by
    \[
        y = e^{-h(t)} \int q(t)\cdot e^{h(t)} \dd{t}
    \]
\end{theoremBox}

\subsection{}

\begin{theoremBox}
    For ODEs of the form
    \[
        \dv{y}{t} + p(t)\cdot y = q(t)
    \]
    the \textbf{general solution} is given by
    \[
        y = v \int \frac{q(t)}{v} \dd{t}
    \]
    where \(v\) satisfies
    \[
        \dv{v}{t} + p(t)\cdot v = 0
    \]
\end{theoremBox}
\begin{proof}
    We want to solve \cref{eq: method3 eq for y}.
    Suppose \(v\) satisfies \cref{eq: method3 eq for v},
    which is separable and can be solved easily.
    \begin{gather}
        y' + p(t)\cdot y = q(t)
        \label{eq: method3 eq for y} \\
        v' + p(t)\cdot v = 0
        \label{eq: method3 eq for v}
    \end{gather}
    Take the derivative of \(y/v\):
    \begin{align*}
        \dv{t}(\frac{y}{v})
        &= \frac{y' v - v' y}{v^2} \\
        &= \frac{y' v + (p(t)\cdot v) y}{v^2} 
        \tag*{Plug in \cref{eq: method3 eq for v}.} \\
        &= \frac{q(t)}{v}
        \tag*{Plug in \cref{eq: method3 eq for y}.}
    \end{align*}
    Thus we have
    \begin{equation*}
        \frac{y}{v} = \int \frac{q(t)}{v} \dd{t}
    \end{equation*}
    Since \(v\) can be solved from \cref{eq: method3 eq for v},
    we obtain the solution for \(y\).
\end{proof}

\section{Non-Linear}

\subsection{Separable Equations}

A first order ODE is \textbf{separable} if it can be written in the form
\[
    M(t) + N(y)\cdot \dv{y}{t} = 0
\]
for some functions \(M\) and \(N\).
Then the solution satisfies
\[
    \int N(y)\dd{y} = -\int M(t)\dd{t}
\]

\subsection{Transformation Methods}

Some complicated non-linear ODEs can be transformed into ODEs that are solvable by suitable substitutions.

\subsubsection{Bernoulli Equations}

Bernoulli equations can be transformed into linear ODEs by the substitution \(v = y^{1-r}\).


\begin{definition}[Bernoulli Equations]
    A first order ODE is called a \textbf{Bernoulli equation} if it can be written in the form
    \begin{equation}
        \dv{y}{t} + p(t)\cdot y = q(t)\cdot y^r
        \label{eq: Bernoulli Equations}
    \end{equation}
    for some given functions \(p\), \(q\) and
    a real number \( r\in\mathbb{R} \setminus \{ 0, 1 \} \).
\end{definition}

\begin{theoremBox}
    The \textbf{general solution} of a Bernoulli equation~\eqref{eq: Bernoulli Equations} is given by
    \[
        y^{1-r} = \frac{1-r}{\mu(t)} \int \mu(t)\cdot q(t)\dd{t}
    \]
    where \(\mu(t)\) satisfies
    \[
        \ln\abs{\mu(t)} = (1-r)\int p(t)\dd{t}
    \]

    Note: If \(r > 0\), then \(y = 0\) is also a solution,
    which is not captured by the above formula.
\end{theoremBox}

\begin{proof}
    Substitute \(v = y^{1-r}\), and then solve with the method of integrating factors.
    \begin{align*}
        \dv{y}{t} + p(t)\cdot y = q(t)\cdot y^r
        \quad & \implies\quad
        y^{-r}\dv{y}{t} + p(t)\cdot y^{1-r} = q(t)
        \qfor y \neq 0 \\
        \quad & \implies\quad
        \dv{t}( y^{1-r}) + (1-r)p(t)\cdot y^{1-r} = (1-r) q(t)
    \end{align*}
\end{proof}

\subsubsection{Homogeneous Equations}

Homogeneous equations can be transformed into separable equations.

\begin{theoremBox}
    If a first-order ODE can be expressed in the following form, then it is said to be \textbf{homogeneous}:
    \[
        \dv{y}{t} = F(y/t)
    \]
    for some given function \(F\).
    
    \noindent
    Then it can be transformed into a separable equation using \(v = y/t\):
    \[
        \frac{1}{F(v) - v} \dd{v} = \frac{1}{t} \dd{t}
    \]
\end{theoremBox}

\begin{remark}
    The homogeneous equations considered here have nothing to do with
    the homogeneous LCCDEs.
\end{remark}

\subsection{Exact Equations}

\begin{definition}[Exact Equations]
    A first order ODE \(M(t,y)\dd{t}+N(t,y)\dd{y}=0\) is said to be \textbf{exact} if there exists a function \(\Psi(t,y)\) such that
    \begin{equation*}
        \pdv{\Psi(t,y)}{t} = M(t,y) \qand \pdv{\Psi(t,y)}{y} = N(t,y)
    \end{equation*}
    Then the general solution is given implicitly as
    \[
        \Psi(t,y) = C \qquad\text{for some constant } C \in \mathbb{R}
    \]
\end{definition}

\begin{theorem}
    Let \(R\) be a simply connected region in the \(t y\)-plane.
    If \(M(t,y)\), \(N(t,y)\), \(M_y(t,y)\), and \(N_t(t,y)\) are continuous in~\(R\), then
    \begin{equation*}
        M(t,y) + N(t,y) \dv{y}{t} = 0 \text{ is exact in } R
        \qquad\iff\qquad
        M_y(t,y) = N_t(t,y) \quad \text{for all } (t,y) \in R
    \end{equation*}
    where \(M_y(t,y) := \pdv*{M(t,y)}{y}\) and \(N_t(t,y) := \pdv*{N(t,y)}{t}\).
\end{theorem}

To solve the first order ODE \(M(t,y)\dd{t}+N(t,y)\dd{y}=0\),
we can follow these steps:
\begin{enumerate}
    \item Verify that the equation is exact by checking if \(M_y(t,y) = N_t(t,y)\).
    \item Integrate \(M(t,y)\) with respect to \(t\) to find \(\Psi(t,y)\):
    \[
        \Psi(t,y) = \int M(t,y)\dd{t} + h(y)
    \]
    where \(h(y)\) is an arbitrary function of \(y\).
    \item Differentiate \(\Psi(t,y)\) with respect to \(y\) and set it equal to \(N(t,y)\):
    \[
        \pdv{\Psi(t,y)}{y} = N(t,y)
    \]
    This will allow us to solve for \(h(y)\).
    \item Substitute \(h(y)\) back into \(\Psi(t,y)\) to obtain the implicit solution:
    \[
        \Psi(t,y) = C
    \]
    for some constant \(C\).
\end{enumerate}

\subsection{Exact Equations and Integrating Factors}

It is sometimes possible to convert an ODE that is not exact into an exact equation with a suitable integrating factor.

Suppose the following first order ODE is exact:
\begin{equation*}
    \mu M(t,y) + \mu N(t,y)\dv{y}{t} = 0
\end{equation*}
Then the integrating factor \( \mu \) must satisfy
\begin{equation*}
    \pdv{(\mu M)}{y} = \pdv{(\mu N)}{t}
\end{equation*}

On one hand, assuming that \( \mu = \mu(t) \), we have
\begin{equation*}
    \dv{\mu}{t} = \frac{M_y - N_t}{N}\cdot \mu
\end{equation*}
If \((M_y - N_t) / N\) depends on \(t\) alone, then we can solve for \(\mu(t)\).

On the other hand, assuming that \( \mu = \mu(y) \), we have
\begin{equation*}
    \dv{\mu}{y} = \frac{N_t - M_y}{M}\cdot \mu
\end{equation*}
If \((N_t - M_y) / M\) depends on \(y\) alone, then we can solve for \(\mu(y)\).

If neither case applies, then try some substitutions to transform the ODE into a solvable form.

\end{document}
